% Shared finite transport on the original probability space.
\subsection{Finite packets at a unit root}
\label{sec:unit-two-step-packets}
Work on the original finite probability space. Put
\[
 A=J+K,\quad T=J-K=J+B,\quad u=B\one=-g,\quad
 p=1-u,\quad q=1+u,\quad d=\|g\|^2,
\]
where $K^*=-K$, $|K(x,y)|\le1$, and $\|K\|^2\le13/32$.
Zero-mass atoms are omitted. Thus $p,q>0$ on the remaining support.
The notation $\bar f_0=\one$ includes the constant vector. Write
$G_n(z)=\sum_{j=0}^{n-1}z^j$ and use the positive-convolution
coefficients $\beta_j$ of \eqref{eq:Q188:conv}.
In particular $\beta_j\ge d-\delta_2\ge\kappa d$ for $j\ge3$,
where $\delta_2=\|K^2\one\|^2$ and $\kappa=19/32$.
Paragraph~\ref{part:Q81.2} supplies the signed-mean bound
\begin{equation}\label{eq:unit-two-step:signed-exponential}
 \E e^{sBW}\le
 e^{s\E BW+(17/64)s^2\|W\|^2}
 \qquad(W\text{ real},\ s\ge0).
\end{equation}

\begin{lemma}[Finite continuation]\label{lem:oc:finite-continuation}
For a real vector $Z$ and an integer $m\ge0$,
\begin{equation}\label{eq:oc:finite-duhamel}
 T^mZ=B^mZ+\sum_{i=0}^{m-1}(\E B^iZ)\bar f_{m-1-i}.
\end{equation}
\end{lemma}
\begin{proof}
The empty sum gives $m=0$. Apply $T$, use
$TB^mZ=B^{m+1}Z+(\E B^mZ)\one$ and
$T\bar f_j=\bar f_{j+1}$, and induct on $m$.
\end{proof}

\begin{lemma}[A finite exponential envelope]
\label{lem:unit-two-step:exponential-envelope}
If $|S|\le b<3$ and $c\ge(1-b/3)^{-1}$, then
$e^S\le1+S+cS^2/2$.
\end{lemma}
\begin{proof}
For $S\ge0$, use $n!\ge2\cdot3^{n-2}$ for $n\ge2$ in the
exponential series. For $S\le0$, its integral remainder is at
most $S^2/2$. Both give the displayed bound.
\end{proof}

\begin{lemma}[The unit-arm logarithm]\label{lem:unit-two-step:unit-log}
For $-1\le u<1$,
\[
 \log(1-u)\le-u-3u^2/10.
\]
\end{lemma}
\begin{proof}
The derivative of $\log(1-u)+u+3u^2/10$ is
$-u(2+3u)/[5(1-u)]$. Its maxima on $[-1,1)$ occur at $-1$
and $0$. Their values are negative and zero, respectively,
since the first four terms of $e^{7/10}$ sum to $12013/6000>2$.
\end{proof}

\begin{theorem}[Unit-root finite-packet transport]
\label{thm:unit-two-step-transport}
\smallskip\noindent
\textit{The positive first row and its continuation.}
Suppose $0\le w\le W$, $Tw=\bar f_r$, $r\ge1$, and
\[
 0<F\le w(s_0+s_1Z+s_2Z^2),\qquad s_0>0,\quad s_2\ge0.
\]
For any real $\lambda$, set
\[
 R=\frac{s_1\lambda}{s_0}Z,
\qquad
 D\ge\frac{s_1}{s_0}\E(wZ)
       +\frac{|s_1|}{s_0}\|w-\lambda\one\|\|Z\|
       +\frac{2Ws_2}{s_0}\|Z\|^2.
\]
For every integer $L\ge1$, with $D_i=\E B^iR$,
\begin{equation}\label{eq:unit-two-step:finite-row}
 0<T^LF\le s_0Y_L,\qquad
 Y_L=\bar f_{r+L-1}+B^LR+D\bar f_{L-1}
             +\sum_{i=1}^{L-1}D_i\bar f_{L-1-i}.
\end{equation}
In particular the scalar $D$ propagates to $D\bar f_{L-1}$.
For $n=r+L-1$, put $m_L=\max(n-3,L-1)$ and, for
$0\le k\le m_L$, define the finite coefficient
\begin{equation}\label{eq:unit-two-step:path-ledger}
 e_k=D\,\mathbf1_{\{k=L-1\}}
 +D_{L-1-k}\,\mathbf1_{\{0\le k\le L-2\}}
 -\beta_{n-k}\,\mathbf1_{\{0\le k\le n-3\}}.
\end{equation}
An indexed term outside its indicated range is absent. For any subset
$\mathcal J\subseteq\{0,\ldots,m_L\}$ the same comparison gives
\begin{equation}\label{eq:unit-two-step:retained-ledger}
 Y_L\le1+G_n(B)u+B^LR
       +\sum_{k\in\mathcal J}e_k\bar f_k
       +2\sum_{k\notin\mathcal J}(e_k)_+.
\end{equation}

\smallskip\noindent
\textit{Pairing with the actual unit arm.}
If $0<G\le Cq+S$, $C>0$, then
\begin{equation}\label{eq:unit-two-step:unit-pairing}
 pG\le C\left(1+\frac{S}{2C}\right)^2,
 \qquad 1+\frac{S}{2C}\ge0.
\end{equation}

\smallskip\noindent
\textit{The complete-vector moment payment.}
Suppose an actual positive message obeys
\[
 0<G\le sY,\quad s>0,\qquad
 Y\le Cq+c+t u+BV,\quad C>0.
\]
For $\nu\in\{4,8\}$,
\begin{equation}\label{eq:unit-two-step:moment-payment}
 \log\E(pG)^\nu\le\nu\log(sC)+\frac\nu C(c+\E BV)
 +\frac1{C^2}\left\{
       \frac{17\nu^2}{64}\|V\|^2
                    +\frac{5\nu}{6}t^2\right\}.
\end{equation}
For one fixed comparison, use the same norm upper bound
$N_0\ge\|V\|^2/C^2$ in both right-hand sides.
Their values satisfy $\mathcal B_4=\mathcal B_8/2-17N_0/4$.

If instead the actual positive message has the square form
\[
 0<H\le\eta(c+BV/2)^2,\qquad c+BV/2\ge0,\quad\eta>0,
\]
then
\begin{equation}\label{eq:unit-two-step:square-log-scale}
 \log\E H^4\le4\log\eta+8(c-1)+4\E BV
                  +\tfrac{17}{4}\|V\|^2.
\end{equation}
The last term may instead be $2\|V\|^2+6\|BV\|^2$.
The underlying exponential estimate, valid for every real $V$, is
\begin{equation}\label{eq:unit-two-step:variance-exponential}
 \E e^{4BV}\le e^{4\E BV+2\|V\|^2+6\|BV\|^2}.
\end{equation}

\end{theorem}

\begin{proof}
Write $wZ=\lambda Z+(w-\lambda)Z$. The identity
$T(wZ)=\E(wZ)+\lambda BZ+B((w-\lambda)Z)$ and the row bound give
\[
|B((w-\lambda)Z)|\le\E|(w-\lambda)Z|
                  \le\|w-\lambda\one\|\|Z\|,
 \qquad T(wZ^2)\le2W\|Z\|^2.
\]
This proves the first positive comparison. Apply the finite identity
\eqref{eq:oc:finite-duhamel} to
$BR$ and then use positivity to obtain
\eqref{eq:unit-two-step:finite-row}. Subtract the exact positive path convolution from $\bar f_n$ in
\eqref{eq:unit-two-step:finite-row}. Its term of index $k$ is
$-\beta_{n-k}\bar f_k$, and the other two terms have exactly the
indices in \eqref{eq:unit-two-step:path-ledger}. On unretained indices use
$0\le\bar f_k\le2$ according to the sign of $e_k$, proving
\eqref{eq:unit-two-step:retained-ledger}.
For \eqref{eq:unit-two-step:unit-pairing}, the upper bound
$Cq+S$ is positive because it dominates $G>0$.
Apply AM--GM to the two nonnegative functions $p$ and
$(Cq+S)/C$, whose sum is $2+S/C$.

For the moment payment use Lemma~\ref{lem:unit-two-step:unit-log},
$\log p\le-u-3u^2/10$ and
$\log Y\le\log C+(Y-C)/C$. Completing the scalar square gives
\[
 \nu t u/C-3\nu u^2/10\le\frac{5\nu}{6C^2}t^2.
\]
The signed-mean exponential estimate at $\nu/C$ now proves
\eqref{eq:unit-two-step:moment-payment}.
Using the same norm upper $N_0$ in both formulas, the other
terms are linear in $\nu$, proving their stated relation.

For \eqref{eq:unit-two-step:square-log-scale}, first keep the
positive factor $\eta$ outside the fourth moment, and use
$\log(c+BV/2)\le c+BV/2-1$. This retains the exact $4\log\eta$.
Here is the variance-sensitive alternative, retaining the same mean.
For a real $W$ with $|W|\le U$, put
$h=\E W^2/U^2\le\theta$ and $m=\E W$. The moment-matching
two-point comparison in paragraph~\ref{part:Q81.2} applies.
If $m\ge0$, its separation is at most $(1+h)U$, so its centered
log moment is at most $s^2U^2(1+3h)/8$, using
$(1+h)^2\le1+3h$. If $m\le0$, write
$w=-m/U$, $a=(h-w^2)/(1+w)^2$. The same Bernoulli estimate gives
the desired coefficient precisely when
\[
 \frac{3(1+a)^4}{4(1+a+a^2)}
 \le\frac{1+3w^2}{(1+w)^2}+3a.
\]
After subtracting $3a$, the left side is
$3(1+a^2)^2/[4(1+a+a^2)]\le3/4$, since
$a^3+a\le\theta^3+\theta<1$; the other side differs from $3/4$
by $(3w-1)^2/[4(1+w)^2]$.
Consequently $\E e^{sW}\le
e^{s\E W+s^2(U^2+3\E W^2)/8}$.
Apply this with $W=BV$, $U=\|V\|$ and $s=4$;
the zero and degenerate moment cases follow directly.

\end{proof}

\begin{corollary}[The original two-edge output]\label{cor:unit-two-step:output}
For an actual nonnegative $H$, put $U=T^2H$ and
$I_q(U)=\E[q^{-1/3}|U|^{8/3}]$. Then
\[
 I_q(U)^3\le\min\left\{
       [\E(U^4/q)]^2,
       e^{(24427/3420)d-5d^2}[\E H^4]^2\right\}.
\]
\end{corollary}
\begin{proof}
H\"older and $\E q=1$ give the first bound. For the second, apply
the generic transfer \eqref{eq:unit-output:ordinary-transfer},
the inverse row \eqref{eq:unit-output:inverse-row}, and its
cost \eqref{eq:unit-output:curvature}. The ordinary norm input
is proved in paragraph~\ref{part:Q64.7.2}.
\end{proof}

\begin{corollary}[The first connector of a square]
\label{cor:unit-two-step:first-square}
Let $m>0$ and use the same original kernel throughout.
If $0<F\le(m+Z/2)^2$, then
\begin{equation}\label{eq:unit-two-step:first-unweighted-square}
 pTF\le(c+BZ/2)^2,\qquad
 c=m+\tfrac12\E Z+\frac{\|Z\|^2}{4m}.
\end{equation}
If instead $0<F\le q(m+Z/2)^2$, $Z=ag-W$, and
$e=\E W$, put
\[
 N=K((m-a)g+W),\qquad
 \Delta\ge m\{\sqrt{d\|Z\|^2}-\E(gZ)\}+\|Z\|^2.
\]
Then
\begin{equation}\label{eq:unit-two-step:first-q-square}
 0<TF\le m^2q+mN-me+\Delta,\qquad
 pTF\le\left(m-e/2+\frac\Delta{2m}+N/2\right)^2.
\end{equation}
Both square bases are nonnegative. The assertions hold for the
actual signed $e,\E(gZ)$ and need no separate sign condition on $q+Z/m$.
\end{corollary}
\begin{proof}
Use the first row of Theorem~\ref{thm:unit-two-step-transport}
with $s_0=m^2$, $s_1=m$, $s_2=1/4$, and $\lambda=1$.
For $w=1$, $r=1$, it gives
$TF\le m^2q+mBZ+m\E Z+\|Z\|^2/2$.
Equation~\eqref{eq:unit-two-step:unit-pairing} gives
\eqref{eq:unit-two-step:first-unweighted-square}.
For $w=q$, $r=2$, use $\bar f_2=q+Kg$ and
$\E Z=-e$ to obtain
\[
 m^2\bar f_2+mBZ=m^2q+mK((m-a)g+W).
\]
The first-row scalar error is at most $-me+\Delta$.
The complete upper bound in \eqref{eq:unit-two-step:first-q-square}
is therefore positive. Apply the same unit pairing with $C=m^2$.
\end{proof}

\begin{corollary}[Bare inward paths]\label{cor:unit-two-step:bare-paths}
For the actual paths on the same measure,
\[
 \E(p\bar f_2)^4\le e^{d/4},\qquad
 \E(p\bar f_3)^4\le e^{-259d/128},\qquad
 \E(p\bar f_n)^4\le e^{-d/2}\quad(n\ge4).
\]
\end{corollary}
\begin{proof}
At $n=2$, $p+\bar f_2=2+Bu$ and $\E Bu=-d$.
For $n\ge3$, retain $\beta_n\bar f_0$ in the positive convolution:
\[
 p+\bar f_n\le2-\beta_n+BG_{n-1}(B)u,
 \qquad \E BG_{n-1}(B)u=-\beta_{n+1}.
\]
The squared norm of $G_2(B)u$ is $d+\delta_2\le45d/32$;
for $n\ge4$, Lemma~\ref{lem:odd-common:spectral} gives
$\|G_{n-1}(B)u\|^2\le d$.
Apply \eqref{eq:unit-two-step:square-log-scale} to the positive
arithmetic-mean square. Its fourth exponent is
$-4\beta_n-4\beta_{n+1}+(17/4)\|G_{n-1}(B)u\|^2$.
At $n=3$, $\beta_3=\beta_4=d$, so this is at most
$-259d/128$; thereafter use
$\beta_n,\beta_{n+1}\ge\kappa d$ and
$8\kappa-17/4=1/2$.
\end{proof}
